% =====================================================================
%  CHƯƠNG 1. CƠ SỞ LÝ THUYẾT MÃ HAMMING (7,4)
%  Nội dung: TV2 (Thanh Trúc), bản workspace ngày 11/11/2026 12:52,
%  đã gồm các chỗ nhóm trưởng sửa (bảng hội chứng, ví dụ r = 1001010,
%  ma trận H cụ thể, bước 3-4 ở 1.8).
%  Số liệu (G, H, bảng 16 từ mã, bảng hội chứng) khớp tài liệu
%  "quy_uoc_chung_dang_he_thong" và đã kiểm tra lại bằng code.
%  Nhãn \label giữ nguyên để Chương 2-4 dẫn chiếu.
%  Câu chuyện Hamming 1947 (thay đoạn "Đặt tình huống" của TV2 theo yêu cầu
%  nhóm trưởng) đã chuyển sang frontmatter/loimodau.tex; ý mã lặp R = 1/3
%  và R = 4/7 nằm ở đó. Chương này mở bằng lộ trình và hộp quy ước.
%  Trình bày (11/11/2026): bố cục, danh sách, bảng, khung định lý/ví dụ theo
%  bộ quy tắc ở config/settings.tex; câu chữ giữ nguyên.
%  Bổ sung theo review 11/11/2026 (14:50): quy ước h_j viết hàng; hộp quy ước
%  nêu đủ các chỗ nhắc dạng cổ điển; phép nhân GF(2); tỉ số 1 lỗi/2 lỗi ở 1.1;
%  bất đẳng thức tam giác và chứng minh Định lý 1.1 ở 1.3; "ê = 0" ở Bước 4.
%  Chỉnh 15/11/2026 theo review lần 3 (claude/review_chuong1-3_lan3):
%  thuật ngữ "phía nhận/phía gửi", "đảo bit", "bit kiểm tra"; "n - k bit dư"
%  (tránh trùng ký hiệu r); bỏ câu lặp cuối khung kết luận 1.1; đoạn chuyển
%  1.1 -> 1.2; GF(2) là trường; |C| = 2^k (1.5, 1.7); ý nghĩa tốc độ mã (1.4);
%  ý nghĩa G H^T = 0 (1.6); Lưu ý 1.3 dẫn chiếu Bảng chế độ giải mã (2.5.2);
%  Hình 1.1 (u -> c -> r -> s -> c^ -> u^) và đoạn tóm tắt cuối 1.8.
%  Chỉnh 15/11/2026 (tối) theo bản nhận xét Chương 1
%  (claude/phan_hoi_nhan_xet_chuong1_2026-11-15):
%  1.1: tỉ số 1 lỗi / 2 lỗi chuyển xuống sau định nghĩa BSC (công thức nhị
%  thức cần giả thiết các bit lỗi độc lập), viết cho khối n bit, giữ hai giá
%  trị p; bảng 4 xác suất điều kiện thay bằng một công thức.
%  1.2: hai cách mô tả không gian con (cơ sở -> G; hệ phương trình -> H) thay
%  cho ma trận G tổng quát (g_ij).
%  1.3: c' -> c~ trong chứng minh Định lý 1.1 (c' chỉ dùng cho từ mã mở rộng).
%  1.4, 1.6: phân biệt mã Hamming với mã tuyến tính (7,4) bất kỳ (phản ví dụ
%  h1 = h5, dmin = 2; tổng quát n = 2^m - 1). 1.5: sửa câu bị ngắt quanh
%  công thức hoán vị. 1.7: rút gọn đoạn mở đầu, nêu bảng dùng lại ở đâu.
% =====================================================================
\chapter{CƠ SỞ LÝ THUYẾT MÃ HAMMING (7,4)}
\label{chap:1}

Chương này đi từ bài toán truyền tin có nhiễu (mục~\ref{sec:1.1}), qua trường
$\GF$ và khoảng cách Hamming (mục~\ref{sec:1.2}--\ref{sec:1.3}), đến việc xây
dựng mã Hamming (7,4) bằng các ma trận $G$, $H$ và sửa lỗi bằng hội chứng
(mục~\ref{sec:1.4}--\ref{sec:1.8}). Trước khi bắt đầu, nhóm thống nhất một số
quy ước dùng xuyên suốt báo cáo.

\begin{quyuoc}[Quy ước dùng trong toàn báo cáo]
Báo cáo dùng mã Hamming (7,4) ở \emph{dạng hệ thống}: từ mã
$c = (u_1u_2u_3u_4\,p_1p_2p_3)$, bốn bit đầu là dữ liệu, ba bit sau là bit kiểm
tra. Các ma trận $G$ và $H$ được xây dựng ở mục~\ref{sec:1.5}
và~\ref{sec:1.6}; dạng cổ điển (bit kiểm tra ở vị trí $1, 2, 4$) được nhắc ở
mục~\ref{sec:1.1}, \ref{sec:1.5}, \ref{sec:1.6} và chỉ được trình bày đầy đủ ở
mục~\ref{sec:2.6}. Các vectơ $u$ (dữ liệu), $c$ (từ mã),
$r$ (vectơ nhận được), $e$ (vectơ lỗi) và $s$ (hội chứng) đều là
\emph{vectơ hàng}; cột thứ $j$ của $H$, ký hiệu $h_j$, cũng được viết như vectơ
hàng (tức $h_j^{\mathsf T}$) khi cộng vào hội chứng. Mã hóa: $c = uG$; hội
chứng: $s = r\HT$. Mọi phép cộng là
cộng modulo~2 (XOR), ký hiệu $\xor$. Các bit được đánh số từ $1$, từ trái sang
phải; trong mã nguồn Python chỉ số bắt đầu từ $0$, nên bit thứ $i$ của báo cáo
là phần tử $i-1$ của danh sách.
\end{quyuoc}

% =====================================================================
\section{Truyền tin có nhiễu và ý tưởng thêm bit dư}
\label{sec:1.1}

\mucnho{1) Truyền tin có nhiễu}
\begin{itemize}
  \item \emph{Khái niệm:} truyền tin có nhiễu là quá trình dịch chuyển thông
        tin từ điểm nguồn đến điểm nhận qua một kênh truyền mà trong đó chịu
        sự tác động của các tín hiệu không mong muốn (gọi là \emph{nhiễu}),
        gây ra hiện tượng sai lệch hoặc suy giảm chất lượng thông tin.
\end{itemize}

\begin{ketluan}
Nhiễu có tính ngẫu nhiên. Phía nhận không thể biết trước khi nào lỗi sẽ xảy
ra; nhiễu chỉ được mô tả qua xác suất lỗi bit $p$ (mô hình ở phần~3). Khi $p$
nhỏ, khối có một lỗi xuất hiện nhiều hơn hẳn khối có hai, ba lỗi, nên một mã
sửa được một lỗi trong mỗi khối đã xử lý được phần lớn các khối bị lỗi.
\end{ketluan}

\mucnho{2) Ý tưởng thêm bit dư}
\begin{itemize}
  \item Khi truyền một chuỗi bit qua đường truyền hoặc lưu vào bộ nhớ, nhiễu
        có thể đảo một số bit từ $0$ thành $1$ hoặc ngược lại. Nếu chỉ gửi
        đúng các bit dữ liệu, phía nhận không có cách nào biết bit nào đã sai.
  \item Mã sửa lỗi giải quyết việc này bằng cách thêm vài bit dư, gọi là
        \emph{bit kiểm tra}. Phía gửi tính các bit kiểm tra từ bit dữ liệu
        theo một quy tắc cố định: dữ liệu gốc $k$ bit được bộ mã hóa gắn thêm
        $n - k$ bit kiểm tra để tạo thành \emph{từ mã} dài $n$ bit.
  \item Vai trò của bit kiểm tra:
        \begin{itemize}
          \item \emph{Phát hiện lỗi:} phía nhận tính lại các bit kiểm tra từ
                dữ liệu nhận được theo cùng quy tắc, rồi so sánh với các bit
                kiểm tra được gửi kèm. Nếu hai kết quả khác nhau, chứng tỏ dữ
                liệu đã bị lỗi trong quá trình truyền.
          \item \emph{Sửa lỗi trực tiếp:} các bit kiểm tra được tính sao cho
                mỗi vị trí lỗi đơn để lại một dấu hiệu (hội chứng) khác nhau.
                Nhờ đó, khi chỉ có một bit bị đảo, phía nhận xác định được bit
                nào sai và chỉ việc đảo lại bit đó để khôi phục dữ liệu gốc.
        \end{itemize}
\end{itemize}

Trong bản gốc, Hamming đặt bit kiểm tra ở các vị trí lũy thừa của 2
($1, 2, 4, \dots$). Báo cáo này dùng \emph{dạng hệ thống}, đưa các bit kiểm
tra về cuối từ mã để thuận tiện cho tính toán ma trận (xem mục~\ref{sec:1.5}
và mục~\ref{sec:2.6}).

\mucnho{3) Giới thiệu về kênh đối xứng nhị phân (BSC) và xác suất lỗi}
\begin{itemize}
  \item \emph{Đặc điểm của kênh đối xứng nhị phân:} mô hình BSC được định
        nghĩa dựa trên ba yếu tố:
        \begin{itemize}
          \item \emph{Tín hiệu vào/ra nhị phân:} kênh truyền chỉ xử lý hai ký
                tự là $0$ và $1$. Đầu vào $X$ và đầu ra $Y$ đều thuộc tập
                $\{0,1\}$.
          \item \emph{Tính chất đối xứng:} xác suất một bit $0$ bị biến đổi
                thành $1$ bằng đúng xác suất một bit $1$ bị biến đổi thành
                $0$.
          \item \emph{Tính không bộ nhớ:} nhiễu tác động lên mỗi bit một cách
                độc lập. Việc bit này bị lỗi không ảnh hưởng đến xác suất lỗi
                của bit tiếp theo.
        \end{itemize}
  \item Gọi $p$ là xác suất lỗi bit, với điều kiện $0 \le p \le 0{,}5$. Với
        cả hai giá trị $x \in \{0, 1\}$ của bit phát, bit nhận được sai với
        xác suất $p$ và đúng với xác suất $1 - p$:
        \[
          \Pr(Y \neq x \mid X = x) = p, \qquad
          \Pr(Y = x \mid X = x) = 1 - p .
        \]
\end{itemize}

\begin{luuy}
Nếu $p > 0{,}5$ thì chỉ cần đảo mọi bit nhận được; khi
đó mỗi bit sai với xác suất $p' = 1 - p < 0{,}5$, tức ta có một kênh BSC tốt
hơn. Vì vậy chỉ cần xét $0 \le p \le 0{,}5$ (BSC sẽ được dùng lại ở
Chương~\ref{chap:5} để tính BER).
\end{luuy}

Vì các bit bị lỗi độc lập với cùng xác suất $p$, xác suất để một khối $n$ bit
có đúng $j$ bit lỗi là $\binom{n}{j} p^j (1-p)^{n-j}$. Với $p > 0$, tỉ số giữa
xác suất có đúng một lỗi và xác suất có đúng hai lỗi là
\[
  \frac{n p (1-p)^{n-1}}{\binom{n}{2} p^2 (1-p)^{n-2}}
  = \frac{2(1-p)}{(n-1)p} .
\]
Với khối $7$ bit, tỉ số này bằng $\dfrac{1-p}{3p}$: khoảng $33$ lần khi
$p = 0{,}01$ nhưng chỉ còn $3$ lần khi $p = 0{,}1$. Đây là cơ sở định lượng cho
kết luận ở phần~1: mã sửa một lỗi trong mỗi khối chỉ hiệu quả khi $p$ nhỏ, và
với cùng $p$, khối càng dài thì tỉ số này càng giảm.

Mỗi khối dữ liệu và mỗi từ mã là một dãy bit $0$, $1$. Để mô tả chính xác quy
tắc tính bit kiểm tra cũng như việc phát hiện và sửa lỗi bằng công cụ của đại
số tuyến tính, mục~\ref{sec:1.2} xem mỗi dãy bit là một vectơ trên trường
$\GF$.

% =====================================================================
\section{Trường \texorpdfstring{$\mathrm{GF}(2)$}{GF(2)} và mã khối tuyến tính \texorpdfstring{$C(n,k)$}{C(n,k)}}
\label{sec:1.2}

\mucnho{1) Trường $\GF$}
\begin{itemize}
  \item \emph{Định nghĩa:} một \emph{trường} là một tập hợp kèm theo hai phép
        toán (cộng và nhân) thỏa mãn các tính chất cấu trúc đại số quen thuộc
        (như tính kết hợp, giao hoán, phân phối, có phần tử nghịch đảo cho
        phép nhân khác không, \dots). Trường Galois 2 phần tử, ký hiệu là
        $\GF$, là trường nhỏ nhất chỉ chứa đúng 2 phần tử: $\{0,1\}$.
  \item Phép cộng trong $\GF$ là phép cộng số học modulo~2, hay phép XOR,
        ký hiệu $\xor$:
\[
  0 \xor 0 = 0,\qquad 0 \xor 1 = 1,\qquad 1 \xor 0 = 1,\qquad 1 \xor 1 = 0 ,
\]
trong đó $1 \xor 1 = 0$ vì trong $\GF$, $1 + 1 = 2 \equiv 0 \pmod 2$. Chính
vì đặc tính này, trong đại số tuyến tính trên trường $\GF$, phép XOR đóng vai
trò là phép cộng vectơ hoặc cộng các thành phần của ma trận.
  \item Phép nhân trong $\GF$ là phép nhân số học thông thường (tương đương
        phép AND):
\[
  0 \cdot 0 = 0,\qquad 0 \cdot 1 = 0,\qquad 1 \cdot 0 = 0,\qquad 1 \cdot 1 = 1 .
\]
Hai bảng trên cho thấy $\GF$ đúng là một trường: phần tử không là $0$, phần
tử đơn vị là $1$; mỗi phần tử là phần tử đối của chính nó ($0 \xor 0 = 0$,
$1 \xor 1 = 0$, nên $-x = x$ với mọi $x$); phần tử khác $0$ duy nhất là $1$
và nghịch đảo nhân của nó là chính nó ($1 \cdot 1 = 1$). Các tính chất kết
hợp, giao hoán và phân phối được thừa hưởng từ phép cộng và phép nhân số
nguyên, vì cả hai phép toán chỉ là tính trên số nguyên rồi lấy số dư khi chia
cho $2$. Tích của hai ma trận trên $\GF$ được tính như thường lệ, trong đó các
tích dùng phép nhân này còn các tổng dùng phép XOR.
  \item Từ trường $\GF$, ta xây dựng không gian vectơ $n$ chiều $\GFn{n}$ gồm
        tất cả các vectơ có độ dài $n$ với các thành phần thuộc $\GF$:
\[
  \GFn{n} = \{\, (x_1, x_2, \dots, x_n) \mid x_i \in \{0,1\} \,\}.
\]
Phép toán trên không gian vectơ này gồm phép cộng vectơ (từng thành phần thực
hiện phép XOR) và phép nhân vô hướng với phần tử trong $\GF$ (nhân với $0$
hoặc $1$).
\end{itemize}

\mucnho{2) Mã khối tuyến tính $C(n,k)$}
\begin{itemize}
  \item Mã khối tuyến tính $C(n,k)$ được định nghĩa là một không gian con
        tuyến tính $C$ có số chiều $k$ của không gian vectơ $\GFn{n}$ (với
        $k < n$).
  \item Vì $C$ là một không gian con của $\GFn{n}$:
        \begin{itemize}
          \item vectơ không $0 \in C$;
          \item tổng của hai từ mã bất kỳ vẫn thuộc $C$ (tính đóng dưới phép
                cộng/XOR);
          \item tích của từ mã với một vô hướng trong $\GF$ vẫn thuộc $C$.
        \end{itemize}
\end{itemize}

\mucnho{3) Vectơ dữ liệu, từ mã và ma trận tạo mã}
\begin{itemize}
  \item \emph{Vectơ dữ liệu:} là chuỗi thông tin gốc cần truyền đi, có độ dài
        là $k$. Ta coi $u$ là một vectơ thuộc không gian vectơ $\GFn{k}$:
        $u = (u_1, u_2, \dots, u_k) \in \GFn{k}$.
  \item \emph{Từ mã:} là chuỗi dữ liệu sau khi đã được mã hóa bằng cách thêm
        các bit kiểm tra lỗi, có độ dài là $n$ ($n > k$). Từ mã $c$ là một
        phần tử thuộc không gian con $C$:
        $c = (c_1, c_2, \dots, c_n) \in C \subset \GFn{n}$.
\end{itemize}

Theo đại số tuyến tính, một không gian con $k$ chiều $C$ của $\GFn{n}$ có thể
được mô tả theo hai cách, và báo cáo dùng cả hai:
\begin{itemize}
  \item \emph{bằng một cơ sở:} $k$ vectơ độc lập tuyến tính sinh ra $C$, xếp
        thành $k$ hàng của \emph{ma trận tạo mã} $G$ cỡ $k \times n$; mỗi từ
        mã là một tổ hợp tuyến tính $c = uG$ của các hàng này
        (mục~\ref{sec:1.5});
  \item \emph{bằng một hệ phương trình tuyến tính thuần nhất:} $n - k$ phương
        trình độc lập mà mọi từ mã phải thỏa, với các hệ số xếp thành $n - k$
        hàng của \emph{ma trận kiểm tra} $H$ cỡ $(n-k) \times n$; khi đó
        $C = \{\, c \in \GFn{n} : c\HT = 0 \,\}$ và $\dim C = n - \rank H$
        (mục~\ref{sec:1.6}; với mã $(7,4)$, chứng minh ở mục~\ref{sec:2.2}).
\end{itemize}
Cách thứ nhất dùng để mã hóa, cách thứ hai dùng để kiểm tra và sửa lỗi; mọi
phép tính trong cả hai cách đều thực hiện trên $\GF$.

% =====================================================================
\section{Khoảng cách Hamming và khả năng phát hiện, sửa lỗi}
\label{sec:1.3}

\mucnho{1) Trọng số Hamming}
\begin{itemize}
  \item \emph{Trọng số Hamming} của một vectơ (hoặc một chuỗi bit) là số
        lượng các phần tử khác không xuất hiện trong vectơ đó.
  \item Đối với trường $\GF$ (chỉ gồm các bit $0$ và $1$), trọng số Hamming của
        một vectơ đơn thuần là tổng số lượng các bit $1$ có trong vectơ đó.
  \item Ký hiệu: trọng số Hamming của vectơ $x$ được ký hiệu là $\wt{x}$.
\end{itemize}

\begin{vidu}
Cho vectơ $x = (1,0,1,1,0)$: có ba bit $1$, vậy trọng số Hamming
$\wt{x} = 3$. Cho vectơ $y = (0,0,0,0)$: không có bit $1$ nào, vậy
$\wt{y} = 0$ (đây là vectơ không).
\end{vidu}

\mucnho{2) Khoảng cách Hamming}

Cho hai vectơ $x = (x_1, x_2, \dots, x_n)$ và $y = (y_1, y_2, \dots, y_n)$ có
cùng độ dài $n$ thuộc $\GFn{n}$. \emph{Khoảng cách Hamming} giữa $x$ và $y$, ký
hiệu $\dist{x}{y}$, được định nghĩa là số lượng vị trí (tọa độ) mà tại đó các
thành phần tương ứng của hai vectơ khác nhau. Vì tại mỗi vị trí, hai bit khác
nhau khi và chỉ khi tổng $\xor$ của chúng bằng $1$, nên
$\dist{x}{y} = \wt{x \xor y}$.

\Needspace{10\baselineskip}
\begin{vidu}
Cho hai chuỗi bit độ dài 7: $x = 1011010$ và $y = 1110010$. So sánh từng vị
trí:
\begin{center}\small
\begin{tabular}{|l|c|c|c|c|c|c|c|}
\hline
\textbf{Vị trí} & 1 & 2 & 3 & 4 & 5 & 6 & 7\\ \hline
$x$ & 1 & 0 & 1 & 1 & 0 & 1 & 0\\
$y$ & 1 & 1 & 1 & 0 & 0 & 1 & 0\\ \hline
So sánh & giống & \textbf{khác} & giống & \textbf{khác} & giống & giống & giống\\ \hline
\end{tabular}
\end{center}
Có tổng cộng 2 vị trí khác nhau (vị trí 2 và vị trí 4), vậy khoảng cách
Hamming $\dist{x}{y} = 2$.
\end{vidu}

\emph{Bất đẳng thức tam giác.} Với mọi $x, y, z \in \GFn{n}$,
$\dist{x}{z} \le \dist{x}{y} + \dist{y}{z}$. Thật vậy,
$x \xor z = (x \xor y) \xor (y \xor z)$ nên
$\dist{x}{z} = \wt{(x \xor y) \xor (y \xor z)}$; mà tại mỗi vị trí mà
$a \xor b$ có bit $1$ thì $a$ hoặc $b$ có bit $1$ tại vị trí đó, nên
$\wt{a \xor b} \le \wt{a} + \wt{b}$.

\mucnho{3) Khả năng phát hiện và sửa chữa lỗi}

Cho một mã khối $C$ (tập hợp tất cả các từ mã hợp lệ). \emph{Khoảng cách tối
thiểu} của mã, ký hiệu $\dmin$, được định nghĩa là khoảng cách Hamming nhỏ nhất
giữa mọi cặp từ mã phân biệt bất kỳ nằm trong $C$:
\[
  \dmin = \min_{c_i \neq c_j} \dist{c_i}{c_j},
  \qquad c_i, c_j \in C .
\]

\begin{dinhly}[Khả năng phát hiện và sửa lỗi]\label{dl:phathien-sua}
Mã có khoảng cách tối thiểu $\dmin$
\begin{enumerate}[label=(\alph*)]
  \item phát hiện được mọi mẫu lỗi có trọng số từ $1$ đến $\dmin - 1$, tức
        phát hiện được tối đa $\dmin - 1$ bit lỗi;
  \item sửa được mọi mẫu lỗi có trọng số không quá
        $t = \left\lfloor \dfrac{\dmin - 1}{2} \right\rfloor$, tức sửa được tối
        đa $t$ bit lỗi.
\end{enumerate}
\end{dinhly}
\begin{proof}

(a) Nếu $r = c \xor e$ với $1 \le \wt{e} \le \dmin - 1$ thì
$r \neq c$ (vì $e \neq 0$). Nếu $r$ cũng là một từ mã thì $r$ và $c$ là hai từ
mã phân biệt với $\dist{c}{r} = \wt{e} < \dmin$, trái với định nghĩa của
$\dmin$. Vậy $r$ không là từ mã và lỗi được phát hiện.

(b) Giả sử $\wt{e} \le t$, tức $\dist{r}{c} \le t$. Vì
$t = \lfloor (\dmin - 1)/2 \rfloor$ nên $2t + 1 \le \dmin$. Với mọi từ mã
$\tilde{c} \neq c$, theo bất đẳng thức tam giác,
$\dist{r}{\tilde{c}} \ge \dist{c}{\tilde{c}} - \dist{r}{c} \ge \dmin - t \ge t + 1
> \dist{r}{c}$. Vậy $c$ là từ mã gần $r$ nhất và là duy nhất, nên bộ giải mã
chọn từ mã gần nhất sẽ khôi phục đúng $c$.
\end{proof}

\mucphu{Ý nghĩa của khả năng phát hiện lỗi (a)}
\begin{itemize}
  \item \emph{Bản chất từ mã:} trong không gian vectơ $\GFn{n}$, tập hợp các
        từ mã hợp lệ $C$ là những điểm ``cô lập'' nằm cách nhau một khoảng
        cách tối thiểu là $\dmin$.
  \item \emph{Điều gì xảy ra khi có lỗi?} Khi một từ mã $c$ được gửi đi, do
        nhiễu kênh truyền, nó bị cộng thêm một vectơ lỗi $e$ (với trọng số
        $1 \le \wt{e} \le \dmin - 1$), biến thành vectơ nhận được
        $r = c \xor e$.
  \item \emph{Tại sao phát hiện được?} Vì số bit lỗi ít hơn khoảng cách tối
        thiểu giữa các từ mã ($\wt{e} < \dmin$), vectơ $r$ nhận được chắc chắn
        không thể rơi vào bất kỳ từ mã hợp lệ nào khác.
  \item \emph{Kết luận:} hệ thống đối chiếu vectơ $r$ với tập hợp từ mã và thấy
        nó ``lơ lửng'', không khớp với bất kỳ từ mã chuẩn nào, nên kết luận:
        ``Đã xảy ra lỗi!'' (dù chưa biết chính xác lỗi ở vị trí nào để sửa).
\end{itemize}

\mucphu{Ý nghĩa của khả năng sửa lỗi (b)}
Hãy tưởng tượng mỗi từ mã
hợp lệ là tâm của một ``quả cầu'' có bán kính là $t$ trong không gian vectơ.
\begin{itemize}
  \item \emph{Nguyên tắc không chồng lấn:} để hệ thống giải mã một cách duy
        nhất (không bị nhầm lẫn), các quả cầu bao quanh các từ mã khác nhau
        tuyệt đối không được đè lên nhau. Khoảng cách giữa tâm của hai từ mã
        phân biệt bất kỳ ít nhất là $\dmin$. Để hai quả cầu bán kính $t$ quanh
        hai từ mã không có điểm chung, cần
        $2t + 1 \le \dmin$, tức $t \le \dfrac{\dmin - 1}{2}$.
  \item \emph{Cơ chế sửa lỗi:} khi phát sinh không quá $t$ bit lỗi, vectơ nhận
        được $r$ sẽ rơi lệch khỏi tâm nhưng vẫn nằm gọn bên trong vùng quả cầu
        của từ mã gốc $c$. Thuật toán giải mã tìm từ mã hợp lệ nằm gần $r$
        nhất (trong cùng quả cầu) và đưa ngược vectơ $r$ về tâm từ mã ban đầu
        $c$.
  \item \emph{Kết luận:} hệ thống không những biết có lỗi mà còn tự động khôi
        phục lại chính xác dữ liệu gốc như chưa từng bị nhiễu.
\end{itemize}

\begin{luuy}
Hai khả năng (a) và (b) ứng với hai chế độ giải mã khác nhau và không cộng dồn
được. Với mã Hamming $(7,4)$ ($\dmin = 3$): ở chế độ chỉ phát hiện, mã phát
hiện mọi mẫu một hoặc hai lỗi; ở chế độ tự động sửa, mã sửa đúng mọi lỗi đơn
nhưng sửa nhầm mọi mẫu hai lỗi (mục~\ref{ss:2.5.3}, Bảng~\ref{tab:chedo}).
\end{luuy}

\begin{menhde}\label{md:dmin-trongso}
Với mã tuyến tính $C$, $\dmin = \min\{\wt{c} : c \in C,\ c \neq 0\}$, tức
$\dmin$ bằng trọng số nhỏ nhất của các từ mã khác $0$.
\end{menhde}
\begin{proof}
Với hai từ mã phân biệt $a, b \in C$ ta có $\dist{a}{b} = \wt{a \xor b}$ và
$a \xor b \in C$, $a \xor b \neq 0$. Ngược lại, mỗi từ mã $c \neq 0$ cho
$\wt{c} = \dist{c}{0}$ với $0 \in C$. Hai tập giá trị trùng nhau nên hai giá
trị nhỏ nhất bằng nhau.
\end{proof}

% =====================================================================
\section{Mã Hamming (7,4)}
\label{sec:1.4}

\begin{itemize}
  \item \emph{Định nghĩa:} mã Hamming $(7,4)$ là một mã khối tuyến tính nhận
        vào một khối $4$ bit dữ liệu thông tin gốc và mã hóa chúng thành một từ
        mã có độ dài $7$ bit bằng cách bổ sung thêm $3$ bit kiểm tra lỗi.
  \item Ký hiệu $C(7,4)$ biểu thị các thông số cốt lõi:
        \begin{itemize}
          \item $n = 7$: độ dài tổng cộng của từ mã (số bit truyền đi trên
                kênh);
          \item $k = 4$: độ dài của vectơ dữ liệu (số bit thông tin thực tế);
          \item $n - k = 3$: số lượng bit kiểm tra;
          \item tốc độ mã $R = \dfrac{k}{n} = \dfrac{4}{7} \approx 0{,}571$
                (tỉ lệ thông tin hữu ích trên tổng dung lượng truyền).
        \end{itemize}
  \item Các thông số trên chưa xác định được mã: có nhiều mã khối tuyến tính
        $(7,4)$ khác nhau, và không phải mã nào cũng sửa được lỗi. Mã Hamming
        là mã mà ba bit kiểm tra được chọn sao cho bảy cột của ma trận kiểm
        tra $H$ là bảy vectơ khác $0$, đôi một khác nhau của $\GFn{3}$
        (mục~\ref{sec:1.6}); nhờ đó mỗi lỗi đơn cho một hội chứng riêng và
        $\dmin = 3$ (mục~\ref{sec:2.3} và~\ref{sec:2.4}).
  \item Tốc độ mã cho biết cái giá của khả năng sửa lỗi: cứ $7$ bit truyền đi
        thì chỉ $4$ bit mang thông tin. Mã lặp ba lần sửa được một lỗi trong
        mỗi khối $3$ bit nhưng chỉ có $R = 1/3$ (xem Lời mở đầu); mã Hamming
        $(7,4)$ cũng sửa được một lỗi trong mỗi khối, với $R = 4/7$. Đổi lại,
        khối dài hơn thì khả năng có từ hai lỗi trong cùng một khối cũng lớn
        hơn (mục~\ref{sec:1.1}).
\end{itemize}

Theo \emph{dạng hệ thống}, từ mã gồm bốn bit dữ liệu đứng trước và ba bit kiểm
tra đứng sau,
\begin{equation}
  c = (c_1\,c_2\,\dots\,c_7) = (u_1\,u_2\,u_3\,u_4\,p_1\,p_2\,p_3),
\end{equation}
trong đó các bit kiểm tra được tính (cộng modulo~2) bởi
\begin{equation}\label{eq:bitkiemtra}
  p_1 = u_1 \xor u_2 \xor u_4,\qquad
  p_2 = u_1 \xor u_3 \xor u_4,\qquad
  p_3 = u_2 \xor u_3 \xor u_4 .
\end{equation}

% =====================================================================
\section{Ma trận tạo mã \texorpdfstring{$G$}{G}}
\label{sec:1.5}

\begin{itemize}
  \item Kích thước: $k \times n$ (với $k$ là số bit dữ liệu, $n$ là độ dài từ
        mã).
  \item Không gian con mã $C$ là một không gian vectơ con $k$ chiều nằm trong
        không gian tổng quát $\GFn{n}$. $k$ hàng của ma trận $G$ là một hệ $k$
        vectơ cơ sở độc lập tuyến tính sinh ra $C$, tức mọi từ mã là tổ hợp
        tuyến tính của các hàng của $G$. Vì các hàng độc lập tuyến tính, hai
        vectơ dữ liệu khác nhau cho hai từ mã khác nhau: nếu $uG = u'G$ thì
        $(u \xor u')G = 0$, tức một tổ hợp tuyến tính các hàng của $G$ với hệ
        số là các bit của $u \xor u'$ bằng $0$, nên $u \xor u' = 0$ và
        $u = u'$. Do đó mã có đúng
        $\lvert C \rvert = \lvert \GFn{k} \rvert = 2^k$ từ mã. Ở dạng chuẩn
        (dạng hệ thống), $G$ viết thành
        \[
          G = [\,I_k \mid P\,],
        \]
        trong đó $I_k$ là ma trận đơn vị kích thước $k \times k$ (giúp giữ
        nguyên vẹn $k$ bit dữ liệu gốc ở đầu từ mã), và $P$ là ma trận phần
        kiểm tra (parity), kích thước $k \times (n-k)$, chứa các quy tắc phối
        hợp bit để sinh ra các bit kiểm tra.
  \item Quá trình mã hóa là một ánh xạ tuyến tính từ không gian $\GFn{k}$ sang
        $\GFn{n}$:
        \begin{equation}\label{eq:mahoa}
          c = u G .
        \end{equation}
  \item Để thuận tiện cho việc trích xuất dữ liệu gốc, các bit được sắp xếp
        sao cho dữ liệu đứng trước, kiểm tra đứng sau. Trong từ mã 7 bit
        $c = (c_1, c_2, \dots, c_7)$:
        \begin{itemize}
          \item vị trí từ $1$ đến $4$ ($c_1, c_2, c_3, c_4$) là $4$ bit dữ liệu
                gốc $u = (u_1, u_2, u_3, u_4)$. Khi nhận được từ mã, ta chỉ
                cần cắt $4$ bit đầu tiên là thu lại ngay dữ liệu ban đầu mà
                không cần giải mã phức tạp;
          \item vị trí từ $5$ đến $7$ ($c_5, c_6, c_7$) là $3$ bit kiểm tra
                ($p_1, p_2, p_3$) được sinh ra thông qua phép nhân ma trận
                $uG$.
        \end{itemize}
\end{itemize}

Cách sắp xếp này tương đương với Hamming cổ điển, chỉ khác nhau ở thứ tự các
cột; để dễ trình bày, báo cáo dùng dạng hệ thống (dữ liệu trước, kiểm tra
sau). Nếu $c_{\text{cổ điển}} = (c_1, \dots, c_7)$ là một từ mã ở dạng cổ
điển thì từ mã ở dạng hệ thống tương ứng là
\begin{equation}\label{eq:hoanvi-1}
  c_{\text{hệ thống}} = (c_3,\,c_5,\,c_6,\,c_7,\,c_1,\,c_2,\,c_4)
\end{equation}
(xem mục~\ref{sec:2.6}).

\begin{vidu}\label{vd:G}
Với mã Hamming $(7,4)$, ma trận phần kiểm tra $P$ có bốn hàng $110$, $101$,
$011$, $111$ và ma trận tạo mã là
\begin{equation}\label{eq:G}
  P = \MatP, \qquad
  G = [\,I_4 \mid P\,] = \MatG .
\end{equation}
Theo dạng hệ thống, $4$ hàng của $G$ đóng vai trò là các vectơ cơ sở, trong đó
$4$ cột đầu là ma trận đơn vị $I_4$ (giữ nguyên dữ liệu) và $3$ cột sau là ma
trận phần kiểm tra $P$ (tạo ra các bit kiểm tra bằng phép XOR). Lập vectơ hàng
$u$ có $4$ bit dữ liệu $u = (1\ 0\ 1\ 1)$ và thực hiện phép nhân
$c = uG$:
\[
  c = (1\ 0\ 1\ 1)\cdot\MatG = (1\ 0\ 1\ 1\ 0\ 1\ 0).
\]
Cách tính thực chất là lấy tổng XOR các hàng của $G$ tương ứng với các vị trí
có số $1$ trong $u$ (các hàng $1$, $3$ và $4$):
\[
  c = \text{Hàng}_1 \xor \text{Hàng}_3 \xor \text{Hàng}_4
    = 1000110 \xor 0010011 \xor 0001111 = 1011010 .
\]
\begin{itemize}
  \item $4$ bit đầu tiên (dữ liệu gốc) chính là $u$, được giữ nguyên nhờ ma
        trận đơn vị $I_4$: $c_{1 \to 4} = (1,0,1,1)$.
  \item $3$ bit sau (bit kiểm tra $p_1, p_2, p_3$) tính theo
        \eqref{eq:bitkiemtra}:
        \begin{align*}
          \text{vị trí 5:}\quad p_1 &= u_1 \xor u_2 \xor u_4 = 1 \xor 0 \xor 1 = 0,\\
          \text{vị trí 6:}\quad p_2 &= u_1 \xor u_3 \xor u_4 = 1 \xor 1 \xor 1 = 1,\\
          \text{vị trí 7:}\quad p_3 &= u_2 \xor u_3 \xor u_4 = 0 \xor 1 \xor 1 = 0 .
        \end{align*}
        Phần kiểm tra là $(0,1,0)$.
\end{itemize}
Vậy từ mã hoàn chỉnh truyền đi là $c = (1,0,1,1,0,1,0)$, trong đó $(1,0,1,1)$
là $4$ bit dữ liệu và $(0,1,0)$ là $3$ bit kiểm tra.
\end{vidu}

% =====================================================================
\section{Ma trận kiểm tra \texorpdfstring{$H$}{H}}
\label{sec:1.6}

\begin{itemize}
  \item Kích thước: $(n-k) \times n$ (với $n-k$ là số lượng bit kiểm tra).
  \item \emph{Bản chất:} nếu các hàng của $G$ sinh ra không gian mã $C$ thì
        $H$ cho các phương trình ràng buộc mà mọi từ mã phải thỏa. Nói cách
        khác, $C$ chính là không gian vô hiệu của $H$:
        \begin{equation}\label{eq:C-ker}
          C = \{\, c \in \GFn{n} : c\HT = 0 \,\}
        \end{equation}
        (chứng minh ở mục~\ref{sec:2.2}).
  \item Nếu ma trận tạo mã có dạng $G = [\,I_k \mid P\,]$ thì ma trận kiểm tra
        tương ứng $H$ có dạng chuẩn
        \[
          H = [\,\PT \mid I_{n-k}\,],
        \]
        với $\PT$ là ma trận chuyển vị của $P$ và $I_{n-k}$ là ma trận đơn vị
        kích thước $(n-k) \times (n-k)$. (Tổng quát $H = [\,-\PT \mid I_{n-k}\,]$;
        trên $\GF$ thì $-\PT = \PT$.)
\end{itemize}

Với mã Hamming $(7,4)$, $P$ có bốn hàng $110$, $101$, $011$, $111$
(mục~\ref{sec:1.5}), nên
\begin{equation}\label{eq:H}
  H = [\,\PT \mid I_3\,] = \MatH .
\end{equation}
Ký hiệu $h_j$ là cột thứ $j$ của $H$, viết dưới dạng hàng ba thành phần (để cộng
được với hội chứng, vốn là vectơ hàng):
\begin{gather*}
  h_1 = \vt{1,1,0},\quad h_2 = \vt{1,0,1},\quad h_3 = \vt{0,1,1},\quad
  h_4 = \vt{1,1,1},\\
  h_5 = \vt{1,0,0},\quad h_6 = \vt{0,1,0},\quad h_7 = \vt{0,0,1}.
\end{gather*}
Bảy cột của $H$ là bảy vectơ khác $0$, đôi một khác nhau của $\GFn{3}$ (lý do
xem mục~\ref{sec:2.3}). Vì $\GFn{3}$ có đúng $2^3 - 1 = 7$ vectơ khác $0$ nên
$H$ dùng hết các vectơ này, mỗi vectơ đúng một lần. Đây là đặc điểm phân biệt
mã Hamming với một mã tuyến tính $(7,4)$ bất kỳ. Chẳng hạn, nếu đổi hàng đầu
của $P$ từ $110$ thành $100$ thì vẫn được một mã tuyến tính $(7,4)$, nhưng khi
đó $h_1 = h_5 = (1,0,0)$: lỗi ở bit $1$ và lỗi ở bit $5$ cho cùng một hội chứng
nên không phân biệt được, và $1000100$ là từ mã trọng số $2$, tức $\dmin = 2$.
Tổng quát, mã Hamming với $m$ bit kiểm tra có ma trận kiểm tra nhận tất cả
$2^m - 1$ vectơ khác $0$ của $\GFn{m}$ làm cột, nên có độ dài $n = 2^m - 1$ và
$k = n - m$ bit dữ liệu; $m = 3$ cho mã $(7,4)$.

\begin{itemize}
  \item \emph{Vai trò trong giải mã và phát hiện lỗi:}
        \begin{itemize}
          \item \emph{Kiểm tra tính hợp lệ:} mọi từ mã chuẩn $c$ khi nhân với
                chuyển vị của $H$ phải triệt tiêu hoàn toàn thành vectơ
                không: $c\HT = 0$.
          \item \emph{Tính hội chứng:} khi nhận được vectơ $r$ (có thể bị
                lỗi), hệ thống tính vectơ $s = r\HT$, gọi là \emph{hội chứng}
                (syndrome). Vì $r\HT$ là tổng các hàng của $\HT$ ứng với các
                vị trí $r_j = 1$, tức tổng các cột $h_j$ của $H$ tại các vị
                trí đó:
                \[
                  s = r\HT = r_1 h_1 \xor r_2 h_2 \xor \dots \xor r_7 h_7 .
                \]
        \end{itemize}
\end{itemize}

Bảng~\ref{tab:hoichung} cho \emph{hội chứng $\to$ vị trí bit lỗi}: bảy cột của
$H$ chính là bảy hội chứng khác nhau, tương ứng với bảy vị trí.

\begin{table}[H]
\centering\small\renewcommand{\arraystretch}{1.1}
\caption{Bảng hội chứng và vị trí bit lỗi (giả thiết tối đa một lỗi)}
\label{tab:hoichung}
\begin{tabular}{|c|c|c|}
\hline
\textbf{Hội chứng} $s = (s_1, s_2, s_3)$ & \textbf{Cột của $H$} &
\textbf{Vị trí bit lỗi}\\ \hline
$(0,0,0)$ & -- & không phát hiện lỗi (không sửa)\\ \hline
$(1,1,0)$ & $h_1$ & vị trí 1\\ \hline
$(1,0,1)$ & $h_2$ & vị trí 2\\ \hline
$(0,1,1)$ & $h_3$ & vị trí 3\\ \hline
$(1,1,1)$ & $h_4$ & vị trí 4\\ \hline
$(1,0,0)$ & $h_5$ & vị trí 5\\ \hline
$(0,1,0)$ & $h_6$ & vị trí 6\\ \hline
$(0,0,1)$ & $h_7$ & vị trí 7\\ \hline
\end{tabular}
\end{table}

Ở dạng hệ thống, hội chứng không đọc ra số thứ tự vị trí mà phải tra bảng;
dạng cổ điển, nơi hội chứng chính là số thứ tự vị trí, được trình bày ở
mục~\ref{sec:2.6}.

\begin{vidu}\label{vd:lat-bit-3}
Đảo bit thứ $3$ của $1011010$ thành $1001010$, tính hội chứng, tìm ra vị trí
$3$ rồi sửa lại. Thay $r = (1,0,0,1,0,1,0)$ vào công thức $s = r\HT$:
\[
  s = 1\cdot h_1 \xor 0\cdot h_2 \xor 0\cdot h_3 \xor 1\cdot h_4
      \xor 0\cdot h_5 \xor 1\cdot h_6 \xor 0\cdot h_7 .
\]
Lấy tổng các cột tại vị trí có bit $1$ trong $r$ (vị trí $1$, $4$, $6$):
\[
  s = h_1 \xor h_4 \xor h_6
    = (1,1,0) \xor (1,1,1) \xor (0,1,0)
    = (1 \xor 1 \xor 0,\ 1 \xor 1 \xor 1,\ 0 \xor 1 \xor 0)
    = (0,1,1).
\]
Vì $s \neq 0$ nên đã phát hiện có lỗi. Tra Bảng~\ref{tab:hoichung}: $(0,1,1)$
trùng cột thứ $3$ của $H$, nên với giả thiết chỉ có một lỗi, bit sai nằm ở vị
trí $3$. Đảo bit tại vị trí thứ $3$ của $r = (1,0,0,1,0,1,0)$ (từ $0$ thành
$1$) được $(1,0,1,1,0,1,0)$. Kết luận: chuỗi từ mã đúng sau khi sửa lỗi là
$1011010$.
\end{vidu}

\mucnho{Mối liên hệ giữa $H$ và $G$}
Dưới góc độ đại số tuyến tính,
tính chất quan trọng nhất liên kết hai ma trận này là \emph{tính trực giao}:
\begin{equation}\label{eq:GHT}
  G\HT = 0 \qquad (\text{hoặc } H G^{\mathsf{T}} = 0),
\end{equation}
được chứng minh ở Mệnh đề~\ref{md:GHT} (mục~\ref{sec:2.2}). Phần tử ở hàng $i$,
cột $j$ của $G\HT$ là tích vô hướng của hàng $i$ của $G$ với hàng $j$ của $H$,
nên hệ thức này nói rằng mỗi hàng của $G$ thỏa cả ba phương trình kiểm tra cho
bởi ba hàng của $H$. Mọi từ mã là tổ hợp tuyến tính các hàng của $G$, nên mọi
từ mã cũng thỏa $c\HT = 0$.

% =====================================================================
\section{Bảng 16 từ mã}
\label{sec:1.7}

Mã Hamming $(7,4)$ có $\lvert C \rvert = 2^4 = 16$ từ mã, mỗi từ mã ứng với
đúng một vectơ dữ liệu $u$ từ $0000$ đến $1111$ (mục~\ref{sec:1.5}). Tính
$c = uG$ cho từng $u$, ta được Bảng~\ref{tab:16tuma}. Cột ``Từ mã'' ghi $4$ bit
dữ liệu, một khoảng trống, rồi $3$ bit kiểm tra.

\begin{table}[H]
\centering\small\renewcommand{\arraystretch}{1.1}
\caption{16 từ mã của mã Hamming (7,4) dạng hệ thống}
\label{tab:16tuma}
\setlength{\tabcolsep}{4pt}
\newcolumntype{C}[1]{>{\centering\arraybackslash}p{#1}}
\begin{tabular}{|C{0.9cm}|C{1.9cm}|C{2.2cm}|C{2.2cm}||C{0.9cm}|C{1.9cm}|C{2.2cm}|C{2.2cm}|}
\hline
\textbf{STT} & \textbf{Vectơ dữ liệu} $u$ & \textbf{Từ mã} $c = uG$ & \textbf{Trọng số Hamming} $\wt{c}$ &
\textbf{STT} & \textbf{Vectơ dữ liệu} $u$ & \textbf{Từ mã} $c = uG$ & \textbf{Trọng số Hamming} $\wt{c}$\\ \hline
1 & 0000 & 0000\,000 & 0 & 9 & 1000 & 1000\,110 & 3\\ \hline
2 & 0001 & 0001\,111 & 4 & 10 & 1001 & 1001\,001 & 3\\ \hline
3 & 0010 & 0010\,011 & 3 & 11 & 1010 & 1010\,101 & 4\\ \hline
4 & 0011 & 0011\,100 & 3 & 12 & 1011 & 1011\,010 & 4\\ \hline
5 & 0100 & 0100\,101 & 3 & 13 & 1100 & 1100\,011 & 4\\ \hline
6 & 0101 & 0101\,010 & 3 & 14 & 1101 & 1101\,100 & 4\\ \hline
7 & 0110 & 0110\,110 & 4 & 15 & 1110 & 1110\,000 & 3\\ \hline
8 & 0111 & 0111\,001 & 4 & 16 & 1111 & 1111\,111 & 7\\ \hline
\end{tabular}
\end{table}

Phân bố trọng số: một từ mã trọng số $0$, bảy từ mã trọng số $3$, bảy từ mã
trọng số $4$ và một từ mã trọng số $7$. Theo Mệnh đề~\ref{md:dmin-trongso},
trọng số nhỏ nhất khác $0$ là $3$, nên $\dmin = 3$ (kiểm tra trực tiếp trên
bảng; chứng minh không cần liệt kê ở mục~\ref{sec:2.4}). Bảng này còn được
dùng để đối chiếu các ví dụ ở Chương~\ref{chap:3} và để suy ra trọng số các từ
mã của mã mở rộng (mục~\ref{ss:2.7.2}).

% =====================================================================
\section{Phát hiện và sửa lỗi bằng hội chứng}
\label{sec:1.8}

\mucnho{Khái niệm}
Khi một từ mã $c$ được truyền đi qua kênh nhiễu, nó
bị tác động bởi một vectơ lỗi (error vector) ký hiệu là $e$. Vectơ nhận được
(ký hiệu là $r$) là tổng của từ mã và vectơ lỗi (thực hiện qua phép XOR):
\[
  r = c \xor e .
\]
Vectơ hội chứng $s$ được định nghĩa là tích giữa vectơ nhận được $r$ và chuyển
vị của ma trận kiểm tra $\HT$:
\[
  s = (s_1\,s_2\,s_3) = r\HT .
\]
Dưới lăng kính đại số tuyến tính, mọi từ mã hợp lệ $c$ đều thỏa mãn phương
trình thuần nhất $c\HT = 0$. Khi phía nhận tính hội chứng $s$ từ vectơ bị nhiễu
$r$, áp dụng tính chất phân phối của phép nhân ma trận:
\begin{equation}\label{eq:hoichung}
  s = r\HT = (c \xor e)\HT = c\HT \xor e\HT = 0 \xor e\HT = e\HT .
\end{equation}
Vậy hội chứng chỉ phụ thuộc vectơ lỗi, không phụ thuộc từ mã đã gửi. Nếu chỉ
có một lỗi ở vị trí $j$ thì $e\HT$ là hàng thứ $j$ của $\HT$, tức $s = h_j$.

Hình~\ref{fig:sodo} tóm tắt đường đi của dữ liệu: mã hóa ở phía gửi, cộng
vectơ lỗi trên kênh, tính hội chứng và sửa lỗi ở phía nhận.

\begin{figure}[H]
\centering
\begin{tikzpicture}[
  node distance=6mm and 9mm,
  khoi/.style={draw,align=center,minimum height=11mm,minimum width=3cm,
               font=\small},
  kenh/.style={khoi,fill=gray!15},
  mui/.style={-{Stealth[length=2.2mm]},thick}]
  \node (u) {$u$};
  \node[khoi,right=6mm of u] (mh) {Mã hóa\\ $c = uG$};
  \node[kenh,right=of mh] (k) {Kênh nhiễu\\ $r = c \xor e$};
  \node[above=4mm of k] (e) {$e$};
  \node[khoi,right=of k] (hc) {Tính hội chứng\\ $s = r\HT$};
  \node[khoi,below=of hc] (sl) {Sửa lỗi\\ $\hat{c} = r \xor \hat{e}$};
  \node[khoi,left=of sl] (gm) {Lấy $4$ bit đầu\\ của $\hat{c}$};
  \node[left=6mm of gm] (uh) {$\hat{u}$};
  \draw[mui] (u) -- (mh);
  \draw[mui] (mh) -- node[above]{$c$} (k);
  \draw[mui] (e) -- (k);
  \draw[mui] (k) -- node[above]{$r$} (hc);
  \draw[mui] (hc) -- node[right]{$s$} (sl);
  \draw[mui] (sl) -- node[above]{$\hat{c}$} (gm);
  \draw[mui] (gm) -- (uh);
\end{tikzpicture}
\caption{Mã hóa, truyền qua kênh nhiễu, tính hội chứng và sửa lỗi}
\label{fig:sodo}
\end{figure}

\mucnho{Quy trình 4 bước phát hiện và sửa lỗi bằng hội chứng}
\begin{enumerate}[label=\textbf{Bước~\arabic*},leftmargin=2.3cm,labelwidth=1.7cm,labelsep=0.6em,align=left]
  \item \textbf{(Nhận dữ liệu):} phía nhận thu được vectơ $r$ từ kênh
        truyền.
  \item \textbf{(Tính hội chứng):} tính $s = r\HT$.
  \item \textbf{(Kiểm tra và phát hiện lỗi):}
        \begin{itemize}
          \item Nếu $s = 0$: $r$ là một từ mã. Bộ giải mã chấp nhận $r$ và lấy
                $4$ bit đầu làm dữ liệu. Kết luận này đúng khi không có lỗi;
                nếu mẫu lỗi trùng một từ mã khác $0$ (cần ít nhất $3$ bit sai)
                thì lỗi không được phát hiện.
          \item Nếu $s \neq 0$: $r$ không phải từ mã, tức đã phát hiện có lỗi.
                Với giả thiết chỉ có đúng một lỗi tại vị trí $i$, vectơ lỗi
                $e$ chỉ có bit $1$ ở vị trí $i$ nên $s = e\HT$ chính là cột $i$
                của $H$; ta đảo bit $i$. Nếu thực tế có từ hai lỗi trở lên thì
                việc đảo bit này luôn cho kết quả sai (bị sửa nhầm thành một từ
                mã khác). Ví dụ: sai bit $1$ và bit $2$ của $1011010$ cho
                $r = 0111010$, $s = h_1 \xor h_2 = (0,1,1) = h_3$, bộ giải mã
                đảo bit $3$ và được $0101010$, là từ mã của $u = 0101$.
        \end{itemize}
  \item \textbf{(Sửa lỗi):} gọi $\hat{e}$ là vectơ chỉ có bit $1$ tại vị
        trí $i$ tìm được ở Bước~3 ($\hat{e} = 0$ nếu $s = 0$). Từ mã sau sửa là $\hat{c} = r \xor \hat{e}$,
        dữ liệu $\hat{u}$ là $4$ bit đầu của $\hat{c}$. Nếu thực tế chỉ có một
        lỗi thì $\hat{e} = e$ và $\hat{c} = c$; nếu có từ hai lỗi trở lên thì
        $\hat{c} \neq c$.
\end{enumerate}

Tóm lại, mã Hamming $(7,4)$ là không gian con $4$ chiều của $\GFn{7}$, sinh bởi
các hàng của $G = [\,I_4 \mid P\,]$ và trùng với không gian nghiệm của
$c\HT = 0$ với $H = [\,\PT \mid I_3\,]$. Bảy cột của $H$ khác $0$ và đôi một
khác nhau nên mỗi lỗi đơn cho một hội chứng riêng, và $\dmin = 3$ nên mã sửa
được một lỗi. Chương~\ref{chap:2} chứng minh các tính chất đã dùng ở đây và mở
rộng sang mã $(8,4)$.
